\documentclass{article}
\usepackage{amsmath}
\usepackage{amssymb}

\begin{document}
\section{Copy paste útiles}

\begin{enumerate}
\item{$\mathbb{Z}$}
\item $x \in l$ 
\item $\wedge$
\item $\vee$
\item $\rightarrow$
\end{enumerate}

\section{Guia 3}

\textbf{Respuestas}\newline

\textbf{Ejercicio 1.}\newline

problema g (n: $\mathbb{Z}$) : $\mathbb{Z}$ \{

    \qquad requiere: \{$n=8 \vee n = 16 \vee n = 131$\}
    
    asegura: \{$(n=8 \rightarrow res = 16) \wedge (n=16 \rightarrow res = 4) \wedge (n=131 \rightarrow res = 1)$\}
\}\newline


\textbf{Ejercicio 2}: Especificar e implementar las siguientes funciones, incluyendo su signatura.\newline


a) absoluto: calcula el valor absoluto de un número entero\newline


problema absoluto ($n: \mathbb{Z}$) : $\mathbb{Z}$ \{

    \qquad requiere: \{True\}

    asegura: \{$res = \sqrt{n^2}$\}
    \}\newline

b) maximoAbsoluto: devuelve el máximo entre el valor absoluto de dos números enteros\newline


problema maximoAbsoluto($n1: \mathbb{Z}, n2: \mathbb{Z}$) : $\mathbb{Z}$ \{

    \qquad requiere: \{True\}

    asegura: \{$res = n1 si absoluto(n1) >= absoluto(n2)$ \}
    asegura: \{$res = n2 si absoluto(n2) > absoluto(n1)$ \}

    \}\newline

c) maximo3: devuelve el máximo entre tres números enteros\newline


problema maximo3($n1: \mathbb{Z}, n2: \mathbb{Z}, n3: \mathbb{Z}$): $\mathbb{Z}$ \{

    \qquad requiere: {True}


    asegura: {$(res = x) \vee (res = y) \vee (res = z)$}
    
    
    asegura: {$(res >= x) \wedge (res >= y) \wedge (res >= z)$}

\}\newline



d) algunoEsCero: dados dos números racionales, decide si alguno es igual a 0 (resolverlo con y sin pattern matching)\newline


problema algunoEsCero($n1: \mathbb{Z}, n2: \mathbb{Z}$) : $\mathbb{Bool}$ \{


    \qquad requiere : \{True\}


    asegura: \{ res = $n1 = 0 \vee n2 = 0$ \}


    \}\newline


e) ambosSonCero: dados dos números reales, decide si ambos son iguales a 0 (resolverlo con y sin pattern matching)\newline


problema ambosSonCero($n1: \mathbb{Z}, n2: \mathbb{Z}$) : $\mathbb{Bool}$ \{
    \qquad requiere : \{True\}
    asegura: \{ res = $n1 = 0 \wedge n2 = 0$ \}
    \}


f) enMismoIntervalo: dados dos números reales, indica si están relacionados por la relación de equivalencia en $\mathbb{R}$ son: (minus inf, 3] , (3,7] y (7,inf), o dicho de otra manera, si pertenecen al mismo intervalo\newline

problema esMismoIntervalo($n1: \mathbb{Z} n2: \mathbb{Z}$): $\mathbb{Bool}$ \{


    \qquad requiere : \{True\}


    asegura: \{res = True si $n1 \in (\infty, 3] \wedge n2 in (\infty, 3] $  \}

    
    asegura: \{res = True si $n1 \in (3, 7] \wedge n2 in (3, 7] $  \}


    asegura: \{res = True si $n1 \in (7, \infty) \wedge n2 in (7, \infty) $  \}

    
    asegura: \{res = False si n1 ni n2 se encuentran en el mismo intervalo\}
    
\}\newline



g) sumaDistintos: que dados tres números naturales, calcule la suma sin sumar repetidos.\newline


problema sumaDistintos($x: \mathbb{R}, y: \mathbb{R}, z: \mathbb{R}$): $\mathbb{R}$ \{

    \qquad requiere:\{True\}


    \qquad asegura: \{si los 3 parámetros son distintos entonces res = x+y+z\}


    \qquad asegura: \{si 2 parámetros son iguales entonces res = al que es distinto\}


    \qquad asegura: \{si los 3 parámetros son iguales entonces res = 0\}

\}\newline



i) digitoUnidades: dado un número entero, extrae su dígito de las unidades.\newline

problema digitoUnidades($x: \mathbb{Z}$): $\mathbb{Z}$ \{

    \qquad requiere: \{True\}
    \qquad asegura: \{res = es el digito de x correspondiente a las unidades\}


\}\newline


j) digitoDecenas: dado un número entero mayor a 9, extrae su dígito de las decenas\newline

problema digitoDecenas($x: \mathbb{Z}$): $\mathbb{Z}$) \{ 
    

    requiere: \{True\}


    asegura: \{ res es el dígito de x correspondiente a las decenas \}

\}\newline

\textbf{Ejercicio 4}: Especificar e implementar las siguientes funciones utilizando tuplas para representar pares y ternas de números.\newline

a) \textbf{productoInterno}: calcula el producto interno entre dos tuplas de $\mathbf{R} x \mathbf{R}$

problema productoInterno($a:\mathbf{R}  x  \mathbf{R} b:\mathbf{R}  x  \mathbf{R}$): ${R}$ \{


    requiere:{True}


    asegura: {$res = a[0]*b[0] + a[1]*b[1]$}

\} \newline

b) \textbf{esParMenor}: dadas dos tuplas de $\mathbf{R}$ x $\mathbf{R}$ , decide si cada coordenada de la primera tupla es menor a la coordenada correspondiente de la segunda tupla. \newline

problema esParMenor($a: \mathbf{R} x \mathbf{R} b: \mathbf{R} x \mathbf{R}$): Bool \{


    requiere: {True}


    asegura: {$res = (a[0]<= b[0]) \wedge (a[1]<=b[1])$}

\} \newline


c) \textbf{distancia}: calcula la distancia euclídea entre dos puntos de $\mathbf{R^2}$ \newline

problema distancia($a: \mathbf{R} x \mathbf{R} b: \mathbf{R} x \mathbf{R}$): $\mathbf{R}$ \{

    requiere:\{True\}

    
    asegura:\{$res = \sqrt{(b[0]-a[0])^2+(b[1]-a[1])^2}$\}


\}\newline


d) \textbf{sumaTerna}: dada una terna de enteros, calcula la suma de sus tres elementos \newline

problema sumaTerna($a: \mathbf{R}x\mathbf{R}x\mathbf{R}$): $\mathbf{R}$ \{
    requiere:\{True\}

    asegura:\{res=a[0]+a[1]+a[2]\}

\} \newline

e) \textbf{sumarSoloMúltiplos}: dada una terna de números enteros y un natural, calcula la suma de los elementos de la terna que son múltiplos del número natural.\newline

problema sumarSoloMúltiplos($a:\mathbf{Z}x\mathbf{Z}x\mathbf{Z}, n:\mathbf{Z}$): $\mathbf{Z}$\{


    requiere: \{n >0\}
    asegura: \{si los tres elementos de la terna son multiplos del natural, res = a[0]+a[1]+a[2] \}
    asegura: \{si un elemento de la terna no es multiplo del natural, res = la suma de los elementos que si son multiplos del natural \}
    asegura:\{si dos elementos de la terna no son multiplos del natural, res = el valor del elemento que sí es multiplo del natural\}
    asegura: \{si ningun elemento de la terna es múltiplo del natural, res = 0\}

\}\newline

f) \textbf{posPrimerPar}: dada una terna de enteros, devuelve la posición del primer número par si es que hay alguno, o devuelve 4 si son todos impares.\newline

PREGUNTAR: ESTÁ BIEN O ES MUY REPETITIVA?

problema posPrimerPar($a: \mathbf{Z}x\mathbf{Z}x\mathbf{Z}$): $\mathbf{Z}$ \{


    requiere: \{True\}


    asegura: \{$res= 0 \vee res =1 \vee res = 2 \vee res =4$\}
    
    
    asegura: \{Si alguno de los elementos es par, entonces el resultado es la posición del elemento que contenga al primer número par\}
    
    
    asegura: \{Si ninguno de los elementos de la terna es par, entonces res = 4\}

\} \newline

g) \textbf{crearPar} a partir de dos componentes, crea un par con esos valores. Debe funcionar para elementos de cualquier tipo \newline

problema crearPar($a: \mathbf{T} b:\mathbf{T}$): $\mathbf{T} x \mathbf{T}$ \{

    requiere: \{True\}


    asegura: \{$(res[0]=a) \wedge (res[1]=b)$\}

\} \newline


h) \textbf{invertir} invierte los elementos del par pasado como parámetro. Debe funcionar para elementos de cualquier tipo. \newline

problema invertir($x: \mathbf{T}x\mathbf{T}$): $\mathbf{T}x\mathbf{T}$ \{


    \qquad requiere: {True}

    \qquad asegura: {$(res[0] = x[1]) \wedge (res[1] = x[0]) $}

\}\newline

\end{document}



